\section*{TD 7 --- Apprentissage non supervisé}
\addcontentsline{toc}{section}{TD 7 --- Apprentissage non supervisé}

\subsection*{Exercice 1 --- Une itération de $k$-means à la main \hfill \normalfont\footnotesize (12 min)}
Quatre points : $P_1=(1,1)$, $P_2=(2,1)$, $P_3=(4,5)$, $P_4=(5,4)$.
Centres initiaux : $C_1=(1,1)$, $C_2=(5,4)$.
\begin{enumerate}
  \item Calculer les distances au carré de chaque point aux deux centres.
  \item Affecter chaque point au centre le plus proche.
  \item Recalculer les deux centres.
  \item Les affectations changeraient-elles à l'itération suivante ?
\end{enumerate}

\subsection*{Exercice 2 --- Choisir $k$ \hfill \normalfont\footnotesize (8 min)}
On trace l'inertie intra-classe en fonction de $k$ : elle décroît de 120 à 40
entre $k=1$ et $k=3$, puis de 40 à 34 entre $k=3$ et $k=8$.
\begin{enumerate}
  \item Où se situe le « coude » ?
  \item Pourquoi l'inertie ne permet-elle pas de choisir $k$ seule ?
  \item Quel autre critère utiliser, et que mesure-t-il ?
\end{enumerate}

\subsection*{Exercice 3 --- ACP \hfill \normalfont\footnotesize (8 min)}
Une ACP sur 4 variables donne des variances expliquées de 60 \%, 25 \%, 10 \%
et 5 \%.
\begin{enumerate}
  \item Combien de composantes faut-il pour dépasser 85 \% ?
  \item Que contient la quatrième composante ?
  \item Pourquoi l'ACP impose-t-elle de standardiser les variables ?
\end{enumerate}

\subsection*{Exercice 4 --- DBSCAN ou $k$-means ? \hfill \normalfont\footnotesize (8 min)}
Indiquer l'algorithme adapté :
\begin{enumerate}
  \item trois groupes de tailles comparables, bien séparés ;
  \item des groupes en forme de croissant ;
  \item des points isolés que l'on veut signaler comme bruit.
\end{enumerate}

\subsection*{Exercice 5 --- Détection d'anomalies \hfill \normalfont\footnotesize (6 min)}
\begin{enumerate}
  \item Pourquoi ne peut-on pas entraîner un classifieur supervisé quand les
        anomalies sont très rares ?
  \item Que fait l'Isolation Forest, en une phrase ?
  \item Quelle métrique utiliser pour évaluer un détecteur d'anomalies ?
\end{enumerate}


\subsection*{Exercice 6 --- Calculer une silhouette \hfill \normalfont\footnotesize (8 min)}
Pour un point : distance moyenne aux points de son groupe $a=2$ ; distance
moyenne au groupe voisin $b=6$.
\begin{enumerate}
  \item Calculer la silhouette $(b-a)/\max(a,b)$.
  \item Que vaut-elle si $a=5$ et $b=5$ ? Comment l'interpréter ?
  \item Quelle est la plage de valeurs possibles ?
\end{enumerate}

\subsection*{Exercice 7 --- Reconstruire après une ACP \hfill \normalfont\footnotesize (7 min)}
On garde une seule composante expliquant 60 \% de la variance.
\begin{enumerate}
  \item Que perd-on en reconstruisant les données d'origine ?
  \item Pourquoi 60 \% peut-il suffire pour visualiser, mais pas pour décider ?
  \item Que mesure la variance expliquée cumulée ?
\end{enumerate}

\subsection*{Exercice 8 --- Minima locaux de $k$-means \hfill \normalfont\footnotesize (7 min)}
Deux initialisations différentes donnent deux partitions différentes.
\begin{enumerate}
  \item Pourquoi l'algorithme dépend-il de l'initialisation ?
  \item Que fait l'initialisation $k$-means++ ?
  \item Comment rendre le résultat plus stable ?
\end{enumerate}

\subsection*{Exercice 9 --- Choisir un seuil d'anomalie \hfill \normalfont\footnotesize (7 min)}
L'Isolation Forest fournit un score continu.
\begin{enumerate}
  \item Pourquoi ne pas prendre un seuil arbitraire à 0,5 ?
  \item Quel coût comparer pour choisir le seuil ?
  \item Quelle métrique évaluer en l'absence d'étiquettes ?
\end{enumerate}

\newpage
